\documentclass[12pt,a4paper]{article}
\usepackage[utf8]{inputenc}
\usepackage{ctex}
\usepackage{amsmath}
\usepackage{amssymb}
\usepackage{graphicx}
\usepackage{hyperref}
\usepackage{geometry}
\usepackage{booktabs}
\geometry{left=2.5cm,right=2.5cm,top=2.5cm,bottom=2.5cm}

\title{量子神经网络的前沿：从参数化量子梳到可扩展的架构搜索}
\author{}
\date{}

\begin{document}

\maketitle

\begin{abstract}
本综述探讨量子神经网络（QNN）实用化面临的可训练性和可设计性双重危机，介绍了三项前沿技术：参数化量子梳（PQComb）、量子神经网络的线性组合（LCQNN）和基于相对波动的量子架构搜索（RF-QAS），为实现可扩展、可训练的QNN设计提供了系统性解决方案。
\end{abstract}

\section{引言——量子机器学习的实用化困境}

\subsection{NISQ时代QNN的机遇与挑战}

在近噪声中等规模量子（NISQ）计算时代，变分量子算法（Variational Quantum Algorithms, VQA）和量子神经网络（Quantum Neural Networks, QNN）被认为是实现量子计算优势最有前途的途径之一。这些混合量子-经典算法利用参数化量子线路（PQC）的高维希尔伯特空间进行计算，并结合经典优化器来调整参数，已在量子化学、组合优化和机器学习等领域展现了巨大潜力。

然而，尽管理论前景广阔，QNN从"原理验证"（proof-of-principle）走向"实用优势"（practical advantage）的道路上却布满了荆棘。其中两个最根本的障碍，构成了当前量子机器学习（QML）领域的实用化困境。

\subsection{QNN的"双重危机"：可训练性与可设计性}

QNN的实用化面临着一对相互交织的根本性挑战：可训练性（Trainability）危机和可设计性（Designability）危机。

\subsubsection{可训练性危机：贫瘠高原（Barren Plateaus, BP）}

可训练性危机的核心是"贫瘠高原"现象。该现象指出，随着量子比特数 $N$ 和电路深度 $D$ 的增加，QNN代价函数的梯度（即参数的偏导数）的方差（variance）会呈指数级 $\mathcal{O}(1/c^N)$ 衰减。这意味着在参数空间的大部分区域，代价函数景观（landscape）都极端平坦，梯度信号微弱到无法被测量和利用。这使得所有基于梯度的优化算法（如随机梯度下降）完全失效，导致QNN在规模扩大后变得无法训练。

\subsubsection{可设计性危机：量子架构搜索（QAS）的扩展性瓶颈}

可设计性危机源于QNN性能对其"架构"（Ansatz，即PQC的线路结构）的极端依赖性。一个精心设计的架构可能是可训练且高性能的，而一个随机或不当的架构则可能毫无用处。然而，手动设计最优架构是一个高度反直觉且几乎不可能完成的任务。自动化"量子架构搜索"（QAS）算法虽然被提出，但其本身面临着严重的扩展性瓶颈：

\begin{itemize}
\item \textbf{搜索空间巨大：}候选架构的数量随 $N$ 和 $D$ 呈指数级增长。
\item \textbf{评估成本高昂：}绝大多数QAS算法需要对每一个候选架构进行完整的训练和验证，才能评估其性能（即"好坏"），这一过程的计算成本极高。
\end{itemize}

这两个危机并非孤立存在，而是深度耦合、互为因果。一个可训练性差（即充满BP）的架构景观会使得QAS的性能评估变得不准确甚至毫无意义，从而误导搜索算法。反之，一个计算成本高昂、无法扩展的QAS算法，即使理论上存在可训练的高性能架构，也永远无法在实际时间内将其找到。因此，要实现QML的真正价值，必须同时攻克这两个瓶颈。

\subsection{本综述的结构与核心}

本综述基于一系列前沿研究成果（主要来自香港科技大学（广州）的研究团队），并结合相关学术文献，旨在深入探讨应对上述"双重危机"的三项关键技术。这些技术共同勾勒出一个从"处理更高阶的量子任务"到"设计可训练、可扩展的QNN"的系统性蓝图。

\begin{itemize}
\item \textbf{参数化量子梳（PQComb）：}本综述将首先介绍PQComb框架。这是一种将QNN的应用范式从传统的"处理量子态"提升到"处理量子过程（如酉变换）"的新型工具，代表了QML能力边界的拓展。
\item \textbf{LCQNN（量子神经网络的线性组合）：}其次，本综述将深入解析LCQNN这一新型QNN架构。该架构旨在正面解决"可训练性危机"，为缓解贫瘠高原问题提供了一种新颖的、可调控的理论框架。
\item \textbf{RF-QAS（基于相对波动的QAS）：}最后，本综述将分析RF-QAS算法。这是一种可扩展的、"免训练"（training-free）的QAS方法，旨在攻克"可设计性危机"中的扩展性瓶颈。
\end{itemize}

本综述将详细阐明这三个框架的理论基础、算法设计、实验性能，并重点分析它们如何协同构成一个解决QNN实用化瓶颈的完整路线图，最后将展望未来的研究方向。

\section{作为量子过程处理器的参数化量子梳（PQComb）}

\subsection{量子过程处理：超越量子态学习}

传统的量子机器学习任务，例如量子态分类，其输入是 $N$ 个量子比特上的一个量子态 $\rho$，输出是一个经典标签 $y$。这是一种 $\text{State} \to \text{Label}$ 的映射。

然而，研究展示了一个更复杂、更强大的任务范式：量子过程处理（Quantum Process Processing）。在"酉逆问题"（unitary inverse problem）中，任务的目标是找到一个量子操作 $C$，使得对于任意输入的未知酉变换 $U$，该操作都能输出其逆操作 $U^{-1}$。即 $C(U^{\otimes N}) = U^{-1}$。在"酉监督学习"中，任务是将一个未知的酉操作 $U$ 作为输入数据，预测其对应的经典标签。

在这两种任务中，输入数据不再是量子态，而是量子操作（酉变换 $U$）本身。这要求计算模型具备一种"量子内存"，能够按顺序接收 $N$ 个 $U$ 的副本，并对这些过程进行相干处理。

\subsection{理论基础：量子梳（Quantum Comb）}

解决此类高阶任务的标准理论工具是量子梳（Quantum Comb）。

\subsubsection{定义与区别}

量子信息论将操作分为不同层级。标准的量子信道（Quantum Channel）（或量子操作）是一个完全正定且保迹（CPTP）的映射，它将一个输入的量子态 $\rho_{in}$ 变换为一个输出的量子态 $\rho_{out}$。

而量子梳（也常被称为超通道（Superchannel）、量子策略（quantum strategy）或过程张量（process tensor））是一种高阶量子操作（Higher-Order Quantum Operation, HOQO）。它接收一个或多个量子操作（如信道 $\mathcal{E}_{in}$）作为输入，并将它们变换为一个新的量子操作 $\mathcal{E}_{out}$。

\subsubsection{数学形式化}

一个 $N$-槽（$N$-slot）量子梳是一个多线性映射。它通过一系列内部的量子操作（"梳齿"，$V_k$）将 $N$ 个输入的量子操作连接起来。为了保证其物理可实现性（即可以构建一个量子线路来实现它），量子梳必须满足严格的因果顺序（causal order）。在数学上，这表现为量子梳的Choi算子 $C$ 必须满足一个"迹条件的层级结构"（hierarchy of trace conditions）。这个层级结构确保了无论在"槽"中插入何种合法的量子操作，整个过程始终是确定性的（即保迹的）。

量子梳是描述和优化量子网络、量子信道辨识、量子计量和量子通信协议的基础框架。

\subsection{PQComb框架（Parameterized Quantum Comb）}

参数化量子梳（PQComb）框架，是将QNN的思想与量子梳理论结合的产物。

\subsubsection{核心思想}

PQComb的核心思想非常直观：将量子梳内部用于连接和处理输入操作的"梳齿"（$V_0, V_1, \ldots, V_m$）参数化。即，使用参数化量子线路（PQC）$V_k(\theta_k)$ 来实现这些内部操作。

\subsubsection{PQComb的结构与优化}

一个典型的PQComb优化任务如下：

\begin{itemize}
\item \textbf{输入（Quantum Dataset）：}一个量子过程数据集 $\{(\mathcal{N}_{j,in}, \mathcal{N}_{j,out})\}_{j=1}^M$。$\mathcal{N}_{j,in}$ 是待处理的输入过程（例如 $U_j^{\otimes N}$），$\mathcal{N}_{j,out}$ 是期望的输出过程（例如 $U_j^{-1}$）。
\item \textbf{模型（Channels for the Comb）：}一个由PQC $\{V_k(\theta_k)\}$ 构成的、参数为 $\theta$ 的量子梳 $C_V(\theta)$。
\item \textbf{优化（Classical Optimization Methods）：}通过混合量子-经典循环，最小化一个损失函数（Loss function），该函数衡量实际输出 $\tilde{\mathcal{N}}_{j,out}(\theta)$ 与期望输出 $\mathcal{N}_{j,out}$ 之间的"不相似度"。
\end{itemize}

优化方法包括：

\paragraph{基于过程的优化 (process-based optimization)}
损失函数 $L_p(\theta)$ 定义为在整个数据集上的平均不相似度：
\begin{equation}
L_p(\theta) = 1 - \sum_j p_j S(\tilde{\mathcal{N}}_{j,out}(\theta), \mathcal{N}_{j,out})
\end{equation}
其中 $S$ 是某种相似度度量（如过程保真度）。

\paragraph{基于梳的优化 (comb-based optimization)}
在特定情况下，如果相似度 $S$ 可以被线性表示，损失函数可以被简化（这在计算上更高效），变为：
\begin{equation}
L_c(\theta) = 1 - \text{Tr}[C_V(\theta)\Omega]
\end{equation}
这里 $C_V(\theta)$ 是PQComb的Choi算子，而 $\Omega$ 是一个固定的、代表目标变换的Choi算子。

\subsubsection{案例分析}

PQComb框架被成功地应用于解决"未知量子比特酉操作反转"问题。通过使用PQComb框架进行优化，研究人员发现了实现该任务的新协议。关键成果是，PQComb找到的协议将实现任意单量子比特酉逆操作所需的辅助（ancilla）量子比特数从先前最佳方法所需的6个减少到了3个，极大地优化了资源需求。

\subsection{PQComb：QML的范式提升}

PQComb框架的意义重大，它不仅是一个精巧的算法，更代表了QML应用的范式提升。它将QNN的能力从"状态处理"（state processing）推广到了"过程处理"（process processing）。

一个 $N$-槽的量子梳在功能上等价于一个"具有内存的量子神经网络"（Quantum neural network with memory）。它能够按顺序接收 $N$ 个"量子程序"（即 $U$）作为输入序列，并对其进行整体运算。这种处理高阶量子动态的能力，为QML在量子控制、量子纠错（例如，实时学习和校正噪声信道）以及量子动力学模拟等需要与量子过程本身交互的领域开辟了全新的应用可能性。

\section{QNN可训练性危机：贫瘠高原（Barren Plateaus）的病理学}

尽管PQComb等框架拓展了QNN的应用边界，但所有QNN（无论处理状态还是过程）都面临着一个共同的、致命的阿喀琉斯之踵：可训练性危机，即贫瘠高原（BP）现象。

\subsection{贫瘠高原（BP）的定义与现象}

BP现象最初由McClean等人在2018年提出，现已成为QNN研究的核心障碍。

BP是指在参数化量子线路（PQC）中，代价函数 $C(\theta)$ 的梯度 $\partial C / \partial \theta_k$ 的方差（Variance）在参数空间中（例如，随机初始化时）的绝大多数区域，都随着量子比特数 $N$ 的增加而呈指数级 $\mathcal{O}(1/c^N)$ 衰减。

梯度幅值在统计上无限接近于零。这意味着代价函数的"景观"（landscape）在几乎所有地方都是平坦的。对于任何基于梯度的优化算法（如SGD, Adam），它们都无法找到有效的下降方向，导致训练完全停滞，QNN变得无法训练。

\subsection{BP的成因分析}

对BP的深入研究揭示了其多种来源：

\begin{itemize}
\item \textbf{全局代价函数（Global Cost Function）：}当代价函数 $C$ 定义为对系统末态的全局测量时（例如，测量 $\langle \psi(\theta) | O | \psi(\theta) \rangle$，其中 $O$ 是一个全局可观测量），BP现象几乎必然出现。这是因为局部的参数变化 $\theta_k$ 对全局测量结果的影响被希尔伯特空间的指数级维度"稀释"了。

\item \textbf{电路深度与随机性（t-Designs）：}随着电路深度 $D$ 的增加，PQC的参数化酉变换会迅速接近一个"酉t-设计"（unitary t-design，通常 $t=2$ 即可）。这意味着PQC在功能上等价于一个从Haar测度随机抽样的酉矩阵。这种高度的随机性导致输出结果高度集中，从而"抹平"了梯度。

\item \textbf{纠缠诱导（Entanglement-induced）BP：}系统中广泛存在的纠缠是BP的另一个重要来源。

\item \textbf{噪声诱导（Noise-induced）BP：}即使是浅层电路，如果使用全局代价函数，量子噪声（如去极化信道）也会指数级地"压缩"梯度方差，导致BP。
\end{itemize}

\subsection{现有的缓解策略及其局限性}

为了克服BP，学术界已经提出了多种缓解策略，它们主要围绕改变代价函数、初始化或架构本身。

\subsubsection{改变代价函数：局域代价函数（Local Cost Functions）}

\textbf{策略：}将全局代价函数分解为多个局域可观测量的总和，例如 $C(\theta) = \sum_i \langle O_i \rangle$，其中 $O_i$ 只作用于少数几个量子比特。

\textbf{效果：}对于足够浅（如 $D \in \mathcal{O}(\log N)$）的电路，这可以避免BP。

\textbf{局限性：}一旦电路深度超过对数尺度，BP依然会重现；此外，许多重要问题（如VQE中的分子哈密顿量）天然需要全局信息，难以完全局域化。

\subsubsection{改变初始化策略}

\textbf{策略：}避免在参数空间中随机初始化。采用特定的初始化方案，如"层级初始化"（layer-wise initialization）、"预训练"（pre-training）或"迁移学习"（transfer learning）。

\textbf{效果：}研究表明，一个好的初始化点（例如从一个可解的局域哈密顿量演化出发）可以为优化提供一个"温暖"的起点，从而避开广阔的BP区域。

\textbf{局限性：}这种方法高度依赖于问题，不能保证全局收敛，且容易陷入局部最优。

\subsubsection{改变架构（Ansatz）}

\textbf{策略：}从根本上限制PQC的结构，使其永远不会形成一个2-design，从而保留梯度信息。

\begin{itemize}
\item \textbf{量子卷积神经网络（QCNN）：}QCNN采用了一种受经典CNN启发的多尺度、分层池化结构。其内在的局域性和层次性被严格证明可以避免BP。

\item \textbf{张量网络（TN）启发的Ansatz：}基于矩阵乘积态（MPS）或树张量网络（qTTN）构建的PQC，其受限的纠缠结构也显示出对BP的抵抗力。

\item \textbf{$k$-局域（$k$-local）Ansatz：}这类PQC仅由 $k$-局域门（即每次最多作用于 $k$ 个量子比特）构成，或者由 $k$-局域哈密顿量演化生成。
\end{itemize}

\subsection{核心权衡：表现力（Expressivity） vs. 可训练性（Trainability）}

上述改变架构的策略是目前最根本的BP缓解方案，但它们引出了QNN设计的核心困境。

\textbf{困境分析：}QCNN、TN-Ansatz或 $k$-局域Ansatz之所以能够保证可训练性，是因为它们主动限制（restrict）了PQC的表现力（Expressivity）。它们无法（也不试图）探索整个指数级的希尔伯特空间。

\textbf{权衡（Trade-off）：}研究证实了这种权衡的普遍存在。
\begin{itemize}
\item 高表现力（例如，深层硬件高效Ansatz）：趋近2-design，导致BP（低可训练性）。
\item 高可训练性（例如，QCNN, 浅层 $k$-local）：表现力受限。
\end{itemize}

\textbf{风险：}如果对架构的限制过于严格（例如，$k$ 值太小），PQC可能根本不具备足够的能力来表示问题的解（例如，一个复杂分子的基态）。在这种情况下，架构虽然"可训练"，但最终会收敛到一个错误的、高能量的解，使其"无用武之地"。

QML领域迫切需要一种新型架构，它能够打破这种"固定的"权衡，提供一种可调控的机制，在表现力和可训练性之间找到针对特定问题的最佳平衡点。

\section{深度解析（一）：LCQNN——平衡表现力与可训练性的新架构}

LCQNN（量子神经网络的线性组合）框架，正是为了解决上述"表现力 vs. 可训练性"的固定权衡问题而提出的。

\subsection{LCQNN（量子神经网络的线性组合）框架}

\subsubsection{设计思想}

LCQNN的核心灵感来自于"酉的线性组合"（Linear Combination of Unitaries, LCU）算法。LCU是现代量子算法（如哈密顿量模拟和HHL算法）的核心子程序。LCU技术通过引入辅助（ancilla）量子比特，并对其进行控制操作，能够相干地实现多个酉操作的线性叠加，即 $\sum_j \alpha_j U_j$。

\subsubsection{LCQNN的架构}

LCQNN架构由两部分量子比特组成：
\begin{itemize}
\item $m$ 个控制（Control）量子比特：处于 $|0\rangle^{\otimes m}$ 初始态。
\item $n$ 个工作（Working）量子比特：用于承载输入数据。
\end{itemize}

其PQC结构如下：

\paragraph{控制层 $V(a)$}
这是一个作用于 $m$ 个控制比特的PQC，由可训练的经典参数 $a$ 控制（例如，$R_y$ 和 $CR_y$ 门）。其目的是将控制比特制备到一个叠加态 $|\psi\rangle = \sum_{j=0}^{L-1} \sqrt{p_j} |j\rangle$，其中 $L \le 2^m$，$p_j$ 是可学习的概率幅。

\paragraph{控制-酉（Control-Unitary）层 $W$}
这是一个大型的（$m+n$）量子比特门，它以 $m$ 个控制比特的状态为索引，对 $n$ 个工作比特执行相应的酉操作：
\begin{equation}
W = \sum_{j=0}^{L-1} |j\rangle\langle j|_m \otimes U_j(\theta_j)_n
\end{equation}

\paragraph{QNN模块 $U_j(\theta_j)$}
这是 $L$ 个独立的、作用于 $n$ 个工作比特的QNN模块（或PQC块），各自拥有独立的参数 $\theta_j$。

\subsubsection{功能}

LCQNN的整体作用是将输入态 $\rho_{in}$ 演化为 $\sum_{j,k} \sqrt{p_j p_k} U_j(\theta_j) \rho_{in} U_k^\dagger(\theta_k)$。通过对辅助比特的测量，该架构可以实现 $L$ 个不同QNN $U_j(\theta_j)$ 的概率性组合或相干叠加。

\subsection{LCQNN的可训练性理论：梯度方差分析}

LCQNN的设计精妙之处在于它如何利用LCU结构来规避贫瘠高原。

\subsubsection{核心洞见}

LCQNN通过将一个庞大、单一的QNN（容易产生BP）分解为 $L$ 个更小、更独立的QNN模块 $U_j(\theta_j)$ 的线性组合，从而将全局的梯度方差"分散"到 $L$ 个独立的模块中，系统性地抑制了BP的产生。

\subsubsection{理论1：梯度方差与 $L$ 的关系}

对于一个由 $L$ 个QNN模块（$L \le 2^m$）组成的LCQNN，其代价函数 $C(\alpha, \theta)$（$\alpha$ 为 $V(a)$ 的参数，$\theta$ 为所有 $U_j$ 的参数）的梯度方差 $\text{Var}[\partial C]$ 满足：
\begin{equation}
\text{Var}_{\alpha, \theta}[\partial C(\alpha, \theta)] \in \mathcal{O}\left(\frac{\text{Poly}(n)}{L^{2n}}\right)
\end{equation}

关键结论是：梯度方差与 $L$（组合数）成反比。这意味着，通过增加 $L$，即增加 $m = \log L$ 个控制比特，可以系统性地、指数级地提升梯度方差，从而消除BP。

\subsubsection{理论2：$k$-局域LCQNN}

当LCQNN中的每个QNN模块 $U_j(\theta_j)$ 被进一步限制为 $k$-局域PQC时（每个门最多作用于 $k$ 个量子比特），可训练性由 $k$ 和 $L$ 共同决定。

更精确的缩放关系（假设 $U_{jt}$ 是 $k$-局域2-design）：
\begin{equation}
\text{Var}_{\alpha, \theta}[\partial C(\alpha, \theta)] \in \mathcal{O}\left(\frac{\text{Poly}(k)}{L^{2k}}\right)
\end{equation}

研究的"Variance Analysis"图清晰地验证了这一理论：
\begin{itemize}
\item 虚线（BP Theory）：标准QNN的梯度方差，随 $n+m$ 呈指数下降（在对数坐标下呈线性下降）。
\item 红线（3-local）/ 橙线（5-local）：$k$-局域LCQNN的梯度方差。它们下降得缓慢得多，成功规避了BP。
\item 关键对比：橙线（5-local）始终位于红线（3-local）之上。这意味着 $k=5$（表现力更强）的架构比 $k=3$（表现力较弱）的架构具有更小的梯度方差。
\end{itemize}

结论：梯度缩放符合理论预期，即局域规模 $k$ 主导可训练性。

\subsection{LCQNN：从"固定权衡"到"可调旋钮"}

LCQNN框架的提出，为解决"表现力 vs. 可训练性"的困境提供了一个全新的视角。

传统的BP缓解策略（如QCNN或 $k$-local Ansatz）提供的是一个固定的权衡点（例如，一个 $k=3$ 的Ansatz）。设计者无法在不破坏其可训练性保证的前提下，轻易地增加其表现力。

而LCQNN提供了一个可调的（tunable）设计。它引入了两个独立的"超参数旋钮"，让设计者可以动态调控：

\begin{itemize}
\item \textbf{旋钮1：$L$（组合数）。}$L$ 由控制比特数 $m$（$L \le 2^m$）决定。根据理论1，增加 $L$ 可以提高可训练性（提升梯度方差）。这实质上开创了一种"以量子比特（辅助比特）换取可训练性"的新型资源权衡策略。

\item \textbf{旋钮2：$k$（局域性）。}$k$ 决定了每个 $U_j$ 模块的内部表现力。根据理论2，增加 $k$（例如从3-local到5-local）可以提高表现力，但这会以降低可训练性（降低梯度方差）为代价。
\end{itemize}

设计者不再被锁定在一个固定的Ansatz上，而是可以根据（1）可用的量子比特总资源（$m+n$）和（2）特定问题的复杂度（需要多大的 $k$），来动态调谐LCQNN架构，在表现力和可训练性之间找到针对该问题的最佳平衡点。

\subsection{性能与应用}

研究证实了LCQNN框架的有效性。该框架在监督学习任务（如MNIST手写数字分类）上进行了验证。实验结果表明，增加 $L$（QNN模块数）和增加每个模块的深度 $D$（即增加 $U_j$ 的表现力），都能够系统性地提高模型的测试准确率。这证明了LCQNN架构的可扩展性及其在真实数据集上学习复杂模式的实用能力。

\section{QNN可设计性危机：量子架构搜索（QAS）的扩展性挑战}

LCQNN虽然提供了一个可训练的"架构模板"（Ansatz class），但它并没有自动回答：对于一个特定问题，最优的 $L$ 和 $k$ 是多少？$U_j(\theta_j)$ 内部的具体门电路应该如何排布？这就引出了QNN的第二个危机：可设计性（Designability）。

\subsection{QAS的必要性与框架}

QNN的性能（包括表现力、可训练性和泛化能力）极度依赖其PQC架构。手动设计一个针对特定问题（如特定分子哈密顿量）的最优架构，是一个极其困难且反直觉的过程。

因此，量子架构搜索（QAS）应运而生。QAS是一种自动化方法，旨在（给定问题、硬件约束等）自动设计一个针对特定性能标准（如最低能量、最高分类精度）优化的PQC架构。

一个QAS框架通常由三个核心组件构成：
\begin{itemize}
\item \textbf{搜索空间（Search Space）：}所有可能的候选架构的集合（例如，门的类型、拓扑结构、$L$ 和 $k$ 的取值）。
\item \textbf{搜索策略（Search Strategy）：}如何在庞大的搜索空间中高效导航（例如，强化学习、遗传算法、梯度下降）。
\item \textbf{评估策略（Evaluation Strategy）：}如何评估一个候选架构的"好坏"（例如，训练后的验证集准确率）。
\end{itemize}

\subsection{主流QAS算法及其核心瓶颈}

目前，主流的QAS搜索策略包括：

\begin{itemize}
\item \textbf{强化学习（RL）-QAS：}将QAS建模为一个马尔可夫决策过程（MDP）。智能体（Agent）通过"添加门"、"选择量子比特"等动作来逐步构建电路，并根据电路的最终性能获得奖励（Reward），从而学习最优的"设计策略"。

\item \textbf{进化算法（EA）-QAS：}模拟生物进化。维护一个PQC架构"种群"，通过"变异"（如随机添加/删除门）、"交叉"（如合并两个架构）和"选择"（保留高性能架构）来迭代优化。

\item \textbf{可微（Differentiable）-QAS (D-QAS)：}借鉴经典NAS的方法，将离散的架构选择（例如，在CNOT门或CZ门之间选择）"松弛"（relax）为连续的概率参数，然后使用梯度下降来同时优化架构参数和门参数。
\end{itemize}

\subsubsection{核心瓶颈：评估策略的高昂成本}

上述所有QAS策略，无论其搜索方式多么智能，都在"评估策略"上遇到了一个共同的、灾难性的瓶颈：为了评估任何一个候选架构的性能，它们都必须在数据集上对该架构进行一次完整的训练和验证。

\textbf{成本分析：}假设搜索空间有 $M$ 个架构，评估（训练）每个架构的平均成本为 $T$。则QAS的总搜索成本为 $\mathcal{O}(M \cdot T)$。

\textbf{灾难性扩展：}$M$（搜索空间）和 $T$（训练成本，受BP影响）本身都是随量子比特数 $N$ 呈指数级增长的。这使得QAS的整体成本呈指数的指数增长，导致QAS在实践中仅限于 $N < 15$ 的小规模问题。

研究的"Searching Time Breakdown"图直观地展示了这一瓶颈。对于传统QAS方法（如DQAS, TF-QAS），其搜索时间随着量子比特数 $N$ 的增加而急剧（近乎指数级）增长。

\section{深度解析（二）：RF-QAS——基于"相对波动"的免训练可扩展搜索}

要攻克QAS的扩展性瓶颈，唯一的出路是切断"评估"与"训练"之间的高成本联系。这催生了"免训练"（Training-Free）QAS的新范式。

\subsection{免训练（Training-Free）QAS：新范式}

\subsubsection{核心思想}

免训练QAS旨在找到一个"零样本"（Zero-Shot）或"免训练"的代理度量（proxy metric）。这个度量 $P(C)$ 必须满足两个条件：
\begin{itemize}
\item \textbf{计算成本低廉：}可以在不训练架构 $C$ 的情况下高效计算出来。
\item \textbf{高度相关：}$P(C)$ 的值与架构 $C$ 经过完整训练后的最终性能（如准确率或能量）高度相关。
\end{itemize}

如果找到这样的 $P(C)$，QAS算法就可以用 $P(C)$ 完全替代昂贵的训练过程，从而将评估成本 $T$ 从指数级降低到多项式级，打破 $\mathcal{O}(M \cdot T)$ 的瓶颈。

\subsection{核心度量：相对波动（RF）作为可学习性的代理}

RF-QAS框架，其核心创新正是找到了这样一个强大的代理度量：RF（相对波动，Relative Fluctuation）。

\subsubsection{RF的定义}

RF被定义为"PQC的问题特定可学习性（problem-specific learnability）的度量"。研究证实了这一新框架的核心是使用"景观波动分析"（landscape fluctuation analysis）。

\subsubsection{RF的工作机制}

RF（也记为 $\tilde{\sigma}$）是一种量化PQC代价函数景观"崎岖程度"或"波动性"的指标。

\begin{itemize}
\item 一个过于平坦的景观（例如BP，RF值极低）显然是不可学习的。
\item 一个过于随机、充满噪声和海量局部极小值的"混沌"景观（RF值可能很高，但缺乏结构）同样难以训练。
\item 一个"好"的、可学习的PQC架构，其代价函数景观应该具有恰到好处的波动性（RF不能太低），以便梯度下降可以跟随其结构；但又不能过于混沌，以确保收敛。
\end{itemize}

RF（相对波动）旨在从统计上捕捉这种"可学习的"景观结构。

\subsubsection{RF的高效计算：Clifford采样}

RF-QAS成功的关键在于其"免训练"特性。它如何能在不训练的情况下评估景观波动性？

研究确认RF度量可以通过"稳定器形式主义"（stabilizer formalism）高效计算。研究中的关键提示："Our QAS method features a lightweight evaluation model enabled by Clifford sampling."

\textbf{机制解析：}RF-QAS的"捷径"在于，它不需要在整个参数空间中进行昂贵的梯度优化来探测景观。相反，它通过在参数空间中进行Clifford采样——即采样那些（在计算上）恰好构成Clifford电路的参数点——来对景观进行统计探测。

由于Clifford电路（由Hadamard, CNOT, S门等构成）的行为可以在经典计算机上通过Gottesman-Knill定理被高效模拟，RF-QAS得以利用高效的经典模拟来统计地估计代价函数景观的"相对波动"（RF）。这使其评估成本 $T$ 变得与训练无关且计算上可行。

\subsection{RF-QAS算法框架}

RF-QAS算法的优势包括："零样本"（Zero-shot）、"可扩展"（scalable）、"改进的准确性和门数量"、"降低的经典资源消耗"。

该算法采用了一种高效的两阶段搜索策略（two-level search strategy）：

\subsubsection{层级搜索（Layer-wise Search）}

算法首先在"层"的级别上进行粗粒度搜索。此阶段的目标是快速确定PQC需要什么样的纠缠模式（例如，哪些量子比特对需要CNOT门，哪些不需要）。

\subsubsection{门级搜索 / 冗余消除（Gate-wise Search / Redundancy Elimination）}

在确定了宏观的层级结构后，算法进入精细的门级优化。此阶段会系统地移除那些对性能（由RF度量评估）影响最小的"冗余"门，以压缩电路，使其更浅、噪声更低。

在这两个阶段中，RF（相对波动）度量被用作唯一的评估标准，指导架构的迭代和剪枝，全程无需任何实际的QNN训练。

\subsection{可扩展性与性能验证}

RF-QAS的有效性得到了令人信服的验证。

\subsubsection{大规模性能（N=50）}

\textbf{任务：}RF-QAS被用于解决大规模（$N=50$）量子多体系统（Ising, Heisenberg, Cluster模型）的基态能量计算（VQE）问题。

\textbf{模拟方法：}研究展示了如何使用MPS（矩阵乘积态）这一经典张量网络方法来高效模拟和评估50比特系统的能量期望值 $\langle \psi | O | \psi \rangle$。

\textbf{结果：}对比了RF-QAS找到的架构与一种强大的、但可能不可扩展的Ansatz（Cartan-2）所能达到的能量。结果显示，RF-QAS找到的架构平均能量差距（Gap）仅为0.01。

\textbf{结论：}这证明了RF-QAS不仅搜索速度快，而且找到的架构是高性能的。

\subsubsection{经典资源消耗（100+ 量子比特）}

\textbf{时间：}"Searching Time Breakdown"图显示，RF-QAS（在图中标记为FluQAS）的搜索时间（Search Time）远低于其他QAS方法（如DQAS, TF-QAS），并且随 $N$ 增长得非常缓慢。

\textbf{内存：}"Memory Usage"图显示，RF-QAS的内存占用增长平缓。即使是100个量子比特的搜索任务，其内存需求（约3.8 GB）也完全在个人电脑的内存框架内。

\textbf{对比：}研究证实，相较于先前局限于6-15个量子比特的QAS方法，RF-QAS在处理50比特问题时，实现了25.23倍的搜索时间缩短和1.41倍的内存改进。

\textbf{总结：}RF-QAS通过其"免训练"的RF度量和高效的Clifford采样评估策略，成功地解决了QAS的扩展性瓶颈，使得在100+量子比特尺度上自动设计QNN架构成为可能。

\section{综合讨论：LCQNN与RF-QAS的协同——QNN设计的未来蓝图}

\subsection{重新审视"双重危机"}

本综述开篇提出了QNN实用化面临的"可训练性"和"可设计性"双重危机。通过前文的深度解析，我们看到：

\begin{itemize}
\item \textbf{LCQNN}通过其LCU架构和（$L, k$）"双旋钮"设计，提供了一个可调控、可保证的"可训练架构模板"，正面解决了"可训练性危机"。
\item \textbf{RF-QAS}通过其"免训练"的RF代理度量和高效采样，提供了一个可扩展的"自动化设计算法"，正面解决了"可设计性危机"。
\end{itemize}

\subsection{LCQNN与RF-QAS的协同作用}

这两项研究虽然是作为独立的工作呈现的，但它们在逻辑上是高度互补的，共同指向了QNN设计的未来蓝图。

\subsubsection{LCQNN的局限}

LCQNN提供了一个"可训练"的架构类别（Search Space）。它通过超参数 $L$（组合数）和 $k$（局域性）来平衡表现力与可训练性。然而，它并没有回答：对于一个特定问题（例如N=50的海森堡模型），$L$ 和 $k$ 应该取多少？$L$ 个 $U_j(\theta_j)$ 模块内部的PQC结构应该是什么样的？手动调优这些新的超参数仍然是一个"可设计性"问题。

\subsubsection{RF-QAS的需求}

RF-QAS提供了一个"可扩展"的搜索方法（Search Algorithm）。它能高效地评估PQC的可学习性（RF度量）。然而，它需要一个定义良好的搜索空间。如果搜索空间是"所有可能的PQC"，这个空间仍然过于庞大且充斥着不可训练的"垃圾"架构。

\subsubsection{协同蓝图：RF-QAS on LCQNN}

一个自然且强大的协同方案是将两者结合：使用RF-QAS算法在LCQNN架构空间中进行搜索。

在这个"大一统"的框架中：
\begin{itemize}
\item \textbf{定义搜索空间：}搜索空间不再是"任意PQC"，而是被限制为"LCQNN类架构"。搜索的（超）参数包括：$L$（组合数）、$k$（局域性），以及每个 $U_j(\theta_j)$ 模块的内部PQC结构（例如门的布局）。
\item \textbf{执行搜索：}运行RF-QAS算法。
\item \textbf{评估：}RF-QAS的"相对波动"（RF）代理度量将自动评估不同（$L, k, \{U_j\}$）组合的"问题特定可学习性"。
\end{itemize}

这种协同框架的巨大优势在于它同时解决了双重危机：
\begin{itemize}
\item \textbf{保证可训练性：}搜索空间（LCQNN）本身具有抑制BP的理论保证。
\item \textbf{保证最优性：}搜索算法（RF-QAS）能高效地在这个空间中找到针对特定问题的最优实例。
\item \textbf{保证可扩展性：}整个搜索过程（RF-QAS）是"免训练"的，因此在计算上是可行的。
\end{itemize}

\subsection{高级量子机器学习框架对比}

为了清晰地总结本综述讨论的三个核心框架的定位和功能，表~\ref{tab:framework_comparison}提供了它们的对比。

\begin{table}[htbp]
\centering
\caption{高级量子机器学习框架对比}
\label{tab:framework_comparison}
\small
\begin{tabular}{p{2cm}p{3cm}p{3cm}p{3cm}p{2.5cm}}
\toprule
\textbf{框架} & \textbf{核心解决的问题} & \textbf{关键创新点} & \textbf{核心度量/理论保障} & \textbf{已验证的可扩展性} \\
\midrule
PQComb & 量子过程处理 & 将QNN从"状态处理器"扩展到"过程处理器" & 基于PQC的变分量子梳 & N/A (重点是功能) \\
\midrule
LCQNN & 可训练性危机 (Barren Plateaus) & LCU架构；以辅助比特换取可训练性 & 梯度方差 $\propto 1/L^{...}$ & N/A (重点是理论) \\
\midrule
RF-QAS & 可设计性危机 (QAS扩展性) & "免训练"的QAS & 相对波动(RF) $\tilde{\sigma}$ & 50-100+ 量子比特 \\
\bottomrule
\end{tabular}
\end{table}

\section{总结与未来研究方向}

\subsection{总结}

本综述从量子神经网络（QNN）面临的"可训练性"（贫瘠高原）和"可设计性"（QAS扩展性）的双重危机出发，系统性地回顾了三个前沿且高度关联的解决方案框架。

\begin{itemize}
\item \textbf{PQComb（参数化量子梳）}将QNN的能力边界从处理"量子态"扩展到了处理高阶的"量子过程"，为量子控制和计量等领域提供了新工具。

\item \textbf{LCQNN（QNN的线性组合）}通过精巧的LCU架构，提供了一种可调控的BP缓解机制。它允许设计者通过调控辅助比特数 $m$（决定 $L$）和局域性 $k$，在表现力和可训练性之间进行动态权衡，成功解决了"可训练性"危机。

\item \textbf{RF-QAS（基于RF的QAS）}通过引入"免训练"的"相对波动"（RF）代理度量，并结合高效的Clifford采样评估，成功解决了QAS的"可扩展性"瓶颈，使得50-100+量子比特规模的QNN自动设计成为可能。
\end{itemize}

这三个框架，特别是LCQNN和RF-QAS的协同，为实现可扩展、可训练且高性能的QNN设计提供了一个清晰的、系统性的蓝图。

\subsection{未来的研究方向}

基于上述分析，未来的研究应集中在以下几个关键方向：

\subsubsection{理论基础深化}

\begin{itemize}
\item \textbf{RF的数学本质：}"相对波动"（RF）目前是作为一种经验性但极其有效的代理度量。未来的研究必须建立RF与PQC可训练性（如梯度方差、收敛速度、Hessian矩阵谱）和泛化能力（Generalization Gap）之间严格的数学和理论联系。

\item \textbf{LCQNN的泛化界：}LCQNN通过增加 $L$（即增加模型参数）来提高可训练性。这是否会增加模型的有效维度或复杂度，从而损害其泛化能力（即过拟合）？对其泛化界的深入理论分析（例如使用PAC-Bayes框架）是评估其最终实用性的关键。
\end{itemize}

\subsubsection{算法与架构的融合}

\textbf{RF-QAS on LCQNN的实现：}立即且最重要的方向是，在算法层面实现两者的融合。即，开发一个以LCQNN架构类（参数为 $L, k, \{U_j\}$）为搜索空间、以RF为评估度量的QAS框架。这将是实现可扩展、可训练的QNN设计的关键一步。

\subsubsection{噪声环境下的鲁棒性}

\begin{itemize}
\item \textbf{噪声感知的RF-QAS：}目前的RF-QAS似乎是在无噪声模拟或轻量级假设下进行的。然而，真实NISQ设备上的噪声本身就是BP的主要来源之一（noise-induced barren plateaus）。未来的研究必须将真实的硬件噪声模型纳入RF度量的计算中，以搜索出在真实设备上具有噪声鲁棒性的QNN架构。

\item \textbf{容错LCQNN：}LCQNN架构依赖于 $m$ 个辅助比特的相干控制（LCU）。这种结构可能对控制比特的退相干噪声高度敏感。研究其在容错量子计算框架下的实现方式和所需的纠错开销，将是评估其长期价值的重要课题。
\end{itemize}

\subsubsection{硬件实验验证}

这三个框架的理论优势最终必须通过真实的硬件实验来验证。

\begin{itemize}
\item 在真实的量子处理器（如超导、离子阱）上执行PQComb的酉逆操作实验。
\item 实验性地测量LCQNN架构的梯度方差随 $L$ 和 $k$ 的变化，以验证其对BP的抑制效果。
\item 在硬件上运行由RF-QAS搜索出的"最优"架构，并与标准的人工设计架构进行性能基准比较。
\end{itemize}

\subsubsection{扩展应用领域}

\begin{itemize}
\item \textbf{PQComb的应用：}将PQComb框架应用于更复杂的量子信息处理任务，例如自适应的量子信道辨识、在线噪声表征（Online Noise Characterization）以及量子纠错码的自适应解码器设计。

\item \textbf{LCQNN/RF-QAS的应用：}将这一套可扩展的QNN设计蓝图，应用于解决当前VQE在量子化学和材料科学中面临的扩展性难题，以及在金融和优化问题中探索真正的量子优势。
\end{itemize}

\section*{参考文献}

\begin{enumerate}
\item McClean, J. R., Boixo, S., Smelyanskiy, V. N., Babbush, R., \& Neven, H. (2018). Barren plateaus in quantum neural network training landscapes. \textit{Nature Communications}, 9(1), 4812.

\item Cerezo, M., Arrasmith, A., Babbush, R., Benjamin, S. C., Endo, S., Fujii, K., ... \& Coles, P. J. (2021). Variational quantum algorithms. \textit{Nature Reviews Physics}, 3(9), 625-644.

\item Bharti, K., Cervera-Lierta, A., Kyaw, T. H., Haug, T., Alperin-Lea, S., Anand, A., ... \& Aspuru-Guzik, A. (2022). Noisy intermediate-scale quantum algorithms. \textit{Reviews of Modern Physics}, 94(1), 015004.

\item Larocca, M., Ju, N., García-Martín, D., Coles, P. J., \& Cerezo, M. (2024). Theory of overparametrization in quantum neural networks. \textit{Nature Computational Science}, 4, 287-297.

\item Chiavazzo, E., Doster, T., Hein, S., Horstmann, B., Levine, Y., Manukian, H., ... \& Zhang, S. X. (2023). Parameterized quantum circuits. \textit{Quantum Science and Technology}, 8(4), 043001.

\item Du, Y., Hsieh, M. H., Liu, T., \& Tao, D. (2020). Expressive power of parametrized quantum circuits. \textit{Physical Review Research}, 2(3), 033125.

\item Huang, H. Y., Broughton, M., Cotler, J., Chen, S., Li, J., Mohseni, M., ... \& McClean, J. R. (2021). Power of data in quantum machine learning. \textit{Nature Communications}, 12(1), 2631.

\item Wang, S., Fontana, E., Cerezo, M., Sharma, K., Sone, A., Cincio, L., \& Coles, P. J. (2021). Noise-induced barren plateaus in variational quantum algorithms. \textit{Nature Communications}, 12(1), 6961.

\item Pesah, A., Cerezo, M., Wang, S., Volkoff, T., Sornborger, A. T., \& Coles, P. J. (2021). Absence of barren plateaus in quantum convolutional neural networks. \textit{Physical Review X}, 11(4), 041011.

\item Sim, S., Johnson, P. D., \& Aspuru-Guzik, A. (2019). Expressibility and entangling capability of parameterized quantum circuits for hybrid quantum-classical algorithms. \textit{Advanced Quantum Technologies}, 2(12), 1900070.

\item Chivilikhin, D., Samarin, A., Ulyantsev, V., Iorsh, I., Oganov, A. R., \& Kyriienko, O. (2020). MoG-VQE: Multiobjective genetic variational quantum eigensolver. arXiv preprint arXiv:2007.04424.

\item Zhang, S. X., Hsieh, C. Y., Zhang, S., \& Yao, H. (2022). Differentiable quantum architecture search. \textit{Quantum Science and Technology}, 7(4), 045023.

\item Chiavazzo, E., \& Hsieh, M. H. (2024). Parameterized quantum comb and simpler circuits for reversing unknown qubit-unitary operations. arXiv preprint arXiv:2403.03761.

\item Qian, Y., Wang, X., Du, Y., Wu, X., \& Tao, D. (2024). LCQNN: Linear combination of quantum neural networks. arXiv preprint arXiv:2507.02832.

\item Chirichella, G., D'Ariano, G. M., \& Perinotti, P. (2009). Theoretical framework for quantum networks. \textit{Physical Review A}, 80(2), 022339.

\item Du, Y., Qian, Y., Wu, X., \& Tao, D. (2025). Scalable quantum architecture search via landscape analysis. arXiv preprint arXiv:2505.05380.

\item Cincio, L., Subaşı, Y., Sornborger, A. T., \& Coles, P. J. (2018). Learning the quantum algorithm for state overlap. \textit{New Journal of Physics}, 20(11), 113022.

\item Havlíček, V., Córcoles, A. D., Temme, K., Harrow, A. W., Kandala, A., Chow, J. M., \& Gambetta, J. M. (2019). Supervised learning with quantum-enhanced feature spaces. \textit{Nature}, 567(7747), 209-212.

\item Chirichella, G., \& Skandera, M. (2022). Quantum combs. In \textit{Quantum Information Processing} (pp. 1-45). Springer.

\item Perinotti, P. (2017). Causal structures and the classification of higher order quantum computations. In \textit{Time in Physics} (pp. 103-127). Springer.

\item Nielsen, M. A., \& Chuang, I. L. (2010). \textit{Quantum Computation and Quantum Information}. Cambridge University Press.
\end{enumerate}

\end{document}
